\subsection{Operational Data Store und Einsatz des KI-Modells zur Kfz-Betrugserkennung}

Wie beschrieben gibt es auch beim Einsatz des KI-Modells zur
Kfz-Betrugserkennung gerichtete Kanten vom Modell-Training 
zu dieser Datenanalyse und von ihr zum Data Warehouse.
Auch hier kommen indirekt die Trainingsdaten zum Einsatz
und die Ergebnisse des Modells werden für zukünftige Analysen persistiert.
Speziell für diesen Anwendungsfall ist,
dass die aktuellen Daten eines Schadensfalls 
für die Analyse, also die Klassifizierung als Betrug oder nicht-Betrug,
in einen Operational Data Store geladen werden.
Für die Historisierung werden sie zusätzlich im Data Warehouse/Data Lake
gespeichert, das ist für diese Analysefragestellung jedoch weniger interessant.

Ein ODS wird hier eingesetzt, weil die Daten schnell für den Zugriff bereit sein sollen.
An den gerichteten Kanten, die zum ODS führen ist notiert, 
dass die Daten hochaktuell und in Batchverarbeitung bereitgestellt werden.
Das ist an dieser Stelle so zu verstehen, 
dass die Daten zu einem Schadensfall gesammelt in den ODS geladen werden,
sobald alle Daten zu diesem Schadensfall in den Datenquellen vorliegen.
Ab diesem Zeitpunkt können die Sachbearbeitenden 
damit beginnen den Fall zu bearbeiten. 
Dabei ist es denkbar, dass ein Schadensfall
noch am Tag, an dem er gemeldet wird, auch bearbeitet wird.
Deshalb wird die höhere Aktualität des ODS benötigt, 
im Vergleich zu den anderen Komponenten dieser Architektur,
die die Daten höchstens tagesaktuell vorhalten.

Gleichzeitig ist der Arbeitsablauf jedoch so ausgelegt,
dass die Sachbearbeitenden die Einschätzung des KI-Modells 
aktiv anfragen, wenn sie einen Schadensfall bearbeiten. 
Es wird also nach dem Pull-Prinzip gearbeitet.
Bei der hohen Aktualität der Daten,
wäre hier auch die Architekturkomponente Real-Time-Processing
denkbar gewesen. 
Allerdings arbeitet das Real-Time-Processing nach dem Push-Prinzip,
und für diesen Anwendungsfall ist es nicht notwendig,
dass die Anwendenden aktiv und automatisiert über die Ergebnisse 
des KI-Modells informiert werden. 

Die Architektur sieht außerdem vor,
dass die Daten über den aktuellen Schadensfall
aus dem ODS gelöscht werden, sobald die Betrugseinschätzung
abgerufen wurde.
So sind zu jedem Zeitpunkt nur die aktuell relevanten Daten
im ODS persistiert und der Speicherbedarf bleibt gering.
Da die Daten zusätzlich in Data Warehouse und Data Lake
historisiert werden, gehen sie dadurch nicht verloren.

Abschließend wird in dem Szenario der Allianz angenommen,
dass die Kfz-Schadenfälle weder für die gesamte Region
zentral noch in jeder Filiale einzeln bearbeitet werden.
Stattdessen gibt es mehrere Schadenszentren pro Region,
die jeweils für einen Teil der Region 
oder für unterschiedliche Schadensarten (Haftpflicht, Kasko, etc.) zuständig sind.
Teilweise werden Fälle auch nach Auslastung der Schadenszentren verteilt. 
Jedes der Schadenszentren erhält einen eigenen ODS, sodass 
personenbezogene Daten nur für die Sachbearbeitenden einsehbar sind, 
die für den Fall zuständig sind.
Daher sind in Abb.~\ref{fig:datenarchitektur-regional} mehrere ODS dargestellt.